% GENERATED FILE. Edit canonical lessons, then run npm run generate:books.
% canonical-source-hash: fnv1a64:28ef5fbecf525712
% canonical-lessons: UR-C392-istaind, UR-C392-istap, UR-C392-rah, UR-C392-tahalna, UR-C392-hanpna

\chapter{Stands and Stops}
\label{ch:ur-stands-and-stops}
\hlchaptermodality{392}

\begin{chapteropening}
\textbf{By the end of this chapter:} \emph{I can find the rickshaw stand and the bus stop, and walk the rest of the way.}

Complete the last lesson of chapter 392: I can find the rickshaw stand and the bus stop, and walk the rest of the way.
\end{chapteropening}

\section[\texorpdfstring{\ur{اسٹینڈ} (isṭainḍ)}{اسٹینڈ (isṭainḍ)}]{\ur{اسٹینڈ} (isṭainḍ) — a stand (for taxis or rickshaws)}

\label{lesson:UR-C392-istaind}



\begin{quote}
Before the new one: say the Urdu for a rush (of people), then the Urdu for to push in.
\end{quote}

\subsection*{You'll want to know: \ur{اسٹینڈ}}
\textbf{\ur{اسٹینڈ}} — \emph{isṭainḍ} — \textquotedblleft{}a stand (for taxis or rickshaws)\textquotedblright{}.

\begin{quote}
\ur{رکشہ} \ur{اسٹینڈ} \ur{بازار} \ur{کے} \ur{سامنے} \ur{ہے۔}
\end{quote}

\begin{quote}
\emph{rikshā isṭainḍ bāzār ke sāmne hai.}
\end{quote}

\begin{quote}
The rickshaw stand is opposite the market.
\end{quote}

Masculine.

\begin{cousinweb}
\textbf{\ur{اسٹینڈ}} (\emph{isṭainḍ}) is borrowed from English \emph{stand}.
\end{cousinweb}

\begin{grammarlens}[title={What you've built}]
One more word for getting around.
\end{grammarlens}

\subsection*{Your turn}
\begin{itemize}
\raggedright
  \item \emph{Say it:} \emph{isṭainḍ}
  \item \emph{Say it:} \emph{isṭainḍ}, once more
  \item \emph{Say it:} \emph{rikshā isṭainḍ bāzār ke sāmne hai.} — The rickshaw stand is opposite the market.
  \item \emph{Recall:} say the Urdu for a rush (of people), then the Urdu for to push in, then say \emph{isṭainḍ} again
\end{itemize}

\begin{tcolorbox}[breakable,colback=teal!4,colframe=teal!35!black,title={Before you move on}]
What does \ur{اسٹینڈ} mean? (\textquotedblleft{}A stand (for taxis or rickshaws)\textquotedblright{}.) Say it once more.
\end{tcolorbox}
\section[\texorpdfstring{\ur{اسٹاپ} (isṭāp)}{اسٹاپ (isṭāp)}]{\ur{اسٹاپ} (isṭāp) — a bus stop}

\label{lesson:UR-C392-istap}



\begin{quote}
Before the new one: say the Urdu for to push in, then the Urdu for a stand (for taxis or rickshaws).
\end{quote}

\subsection*{You'll want to know: \ur{اسٹاپ}}
\textbf{\ur{اسٹاپ}} — \emph{isṭāp} — \textquotedblleft{}a bus stop\textquotedblright{}.

\begin{quote}
\ur{اگلا} \ur{اسٹاپ} \ur{اسکول} \ur{کے} \ur{پاس} \ur{ہے۔}
\end{quote}

\begin{quote}
\emph{aglā isṭāp iskūl ke pās hai.}
\end{quote}

\begin{quote}
The next stop is near the school.
\end{quote}

Masculine.

\begin{cousinweb}
\textbf{\ur{اسٹاپ}} (\emph{isṭāp}) is borrowed from English \emph{stop}.
\end{cousinweb}

\begin{grammarlens}[title={What you've built}]
One more word for getting around.
\end{grammarlens}

\subsection*{Your turn}
\begin{itemize}
\raggedright
  \item \emph{Say it:} \emph{isṭāp}
  \item \emph{Say it:} \emph{isṭāp}, once more
  \item \emph{Say it:} \emph{aglā isṭāp iskūl ke pās hai.} — The next stop is near the school.
  \item \emph{Recall:} say the Urdu for to push in, then the Urdu for a stand (for taxis or rickshaws), then say \emph{isṭāp} again
\end{itemize}

\begin{tcolorbox}[breakable,colback=teal!4,colframe=teal!35!black,title={Before you move on}]
What does \ur{اسٹاپ} mean? (\textquotedblleft{}A bus stop\textquotedblright{}.) Say it once more.
\end{tcolorbox}
\section[\texorpdfstring{\ur{راہ} (rāh)}{راہ (rāh)}]{\ur{راہ} (rāh) — a path; a way}

\label{lesson:UR-C392-rah}



\begin{quote}
Before the new one: say the Urdu for a stand (for taxis or rickshaws), then the Urdu for a bus stop.
\end{quote}

\subsection*{You'll want to know: \ur{راہ}}
\textbf{\ur{راہ}} — \emph{rāh} — \textquotedblleft{}a path; a way\textquotedblright{}.

\begin{quote}
\ur{یہ} \ur{راہ} \ur{سیدھی} \ur{دریا} \ur{تک} \ur{جاتی} \ur{ہے۔}
\end{quote}

\begin{quote}
\emph{yih rāh sīdhī daryā tak jātī hai.}
\end{quote}

\begin{quote}
This path runs straight to the river.
\end{quote}

Feminine.

\begin{cousinweb}
\textbf{\ur{راہ}} (\emph{rāh}) is usually traced to Persian \emph{rāh}, \textquotedblleft{}a road, a way\textquotedblright{}; it is the first half of \textbf{\ur{راہگیر}}, \textquotedblleft{}a passer-by\textquotedblright{}.
\end{cousinweb}

\begin{grammarlens}[title={What you've built}]
One more word for getting around.
\end{grammarlens}

\subsection*{Your turn}
\begin{itemize}
\raggedright
  \item \emph{Say it:} \emph{rāh}
  \item \emph{Say it:} \emph{rāh}, once more
  \item \emph{Say it:} \emph{yih rāh sīdhī daryā tak jātī hai.} — This path runs straight to the river.
  \item \emph{Recall:} say the Urdu for a stand (for taxis or rickshaws), then the Urdu for a bus stop, then say \emph{rāh} again
\end{itemize}

\begin{tcolorbox}[breakable,colback=teal!4,colframe=teal!35!black,title={Before you move on}]
What does \ur{راہ} mean? (\textquotedblleft{}A path; a way\textquotedblright{}.) Say it once more.
\end{tcolorbox}
\section[\texorpdfstring{\ur{ٹہلنا} (ṭahalnā)}{ٹہلنا (ṭahalnā)}]{\ur{ٹہلنا} (ṭahalnā) — to stroll; to walk up and down}

\label{lesson:UR-C392-tahalna}



\begin{quote}
Before the new one: say the Urdu for a bus stop, then the Urdu for a path.
\end{quote}

\subsection*{You'll want to know: \ur{ٹہلنا}}
\textbf{\ur{ٹہلنا}} — \emph{ṭahalnā} — \textquotedblleft{}to stroll; to walk up and down\textquotedblright{}.

\begin{quote}
\ur{شام} \ur{کو} \ur{ہم} \ur{دریا} \ur{کے} \ur{کنارے} \ur{ٹہلتے} \ur{ہیں۔}
\end{quote}

\begin{quote}
\emph{shām ko ham daryā ke kināre ṭahalte haiṅ.}
\end{quote}

\begin{quote}
In the evening we stroll along the riverbank.
\end{quote}

Past (he …): \ur{وہ} \ur{پارک} \ur{میں} \ur{ٹہلا}, \emph{voh pārk meṅ ṭahlā}.

\begin{cousinweb}
\textbf{\ur{ٹہلنا}} (\emph{ṭahalnā}) is inherited: CDIAL lists the Hindi cousin \emph{ṭahalnā}, \textquotedblleft{}to take a walk\textquotedblright{}, under \emph{ṭahall}, a form it reconstructs (CDIAL 5453).
\end{cousinweb}

\begin{grammarlens}[title={What you've built}]
One more word for getting around.
\end{grammarlens}

\subsection*{Your turn}
\begin{itemize}
\raggedright
  \item \emph{Say it:} \emph{ṭahalnā}
  \item \emph{Say it:} \emph{ṭahalnā}, once more
  \item \emph{Say it:} \emph{shām ko ham daryā ke kināre ṭahalte haiṅ.} — In the evening we stroll along the riverbank.
  \item \emph{Recall:} say the Urdu for a bus stop, then the Urdu for a path, then say \emph{ṭahalnā} again
\end{itemize}

\begin{tcolorbox}[breakable,colback=teal!4,colframe=teal!35!black,title={Before you move on}]
What does \ur{ٹہلنا} mean? (\textquotedblleft{}To stroll; to walk up and down\textquotedblright{}.) Say it once more.
\end{tcolorbox}
\section[\texorpdfstring{\ur{ہانپنا} (hāṅpnā)}{ہانپنا (hāṅpnā)}]{\ur{ہانپنا} (hāṅpnā) — to pant}

\label{lesson:UR-C392-hanpna}



\begin{quote}
Before the new one: say the Urdu for a path, then the Urdu for to stroll.
\end{quote}

\subsection*{You'll want to know: \ur{ہانپنا}}
\textbf{\ur{ہانپنا}} — \emph{hāṅpnā} — \textquotedblleft{}to pant\textquotedblright{}.

\begin{quote}
\ur{سارا} \ur{دوڑ} \ur{کر} \ur{اسٹاپ} \ur{پر} \ur{پہنچی} \ur{اور} \ur{ہانپ} \ur{رہی} \ur{تھی۔}
\end{quote}

\begin{quote}
\emph{sārā dauṛ kar isṭāp par pahuṅchī aur hāṅp rahī thī.}
\end{quote}

\begin{quote}
Sara ran to the stop and was out of breath.
\end{quote}

Past (he …): \ur{وہ} \ur{سیڑھی} \ur{پر} \ur{ہانپا}, \emph{voh sīṛhī par hāṅpā}.

\begin{cousinweb}
\textbf{\ur{ہانپنا}} (\emph{hāṅpnā}) is inherited from \emph{hamph}, an imitative form CDIAL reconstructs for panting (CDIAL 13973).
\end{cousinweb}

\begin{grammarlens}[title={What you've built}]
One more word for getting around.
\end{grammarlens}

\subsection*{Your turn}
\begin{itemize}
\raggedright
  \item \emph{Say it:} \emph{hāṅpnā}
  \item \emph{Say it:} \emph{hāṅpnā}, once more
  \item \emph{Say it:} \emph{sārā dauṛ kar isṭāp par pahuṅchī aur hāṅp rahī thī.} — Sara ran to the stop and was out of breath.
  \item \emph{Recall:} say the Urdu for a path, then the Urdu for to stroll, then say \emph{hāṅpnā} again
\end{itemize}

\begin{tcolorbox}[breakable,colback=teal!4,colframe=teal!35!black,title={Before you move on}]
What does \ur{ہانپنا} mean? (\textquotedblleft{}To pant\textquotedblright{}.) Say it once more.
\end{tcolorbox}
